\section{Methodology}
\label{sec:methodology}

In this section, we present the mathematical formulation of our proposed framework, \textbf{PolyMerge}. We first review the standard formulation of task vectors and adaptive model merging, define the Overfitting-Optimizer Paradox mathematically, and then detail the polynomial parameterization and optimization process of PolyMerge.

\subsection{Preliminaries and Setup}
Let $\Theta_{\text{base}} \in \mathbb{R}^M$ represent the weights of a pre-trained base model (e.g., a CLIP ViT-B/32 backbone). Suppose we are given $K$ task-specific expert models, each fine-tuned from the same base model on a distinct downstream task $k \in \{1, 2, \dots, K\}$, resulting in weights $\Theta_k \in \mathbb{R}^M$. Following \cite{ilharco2022editing}, we define the task vector $\mathbf{\Delta}_k \in \mathbb{R}^M$ for each task $k$ as the difference between the expert weights and the base weights:
\begin{equation}
\mathbf{\Delta}_k = \Theta_k - \Theta_{\text{base}}
\end{equation}
To allow for structural variation across depth, we partition the model's weights into $L$ sequential layer-specific parameter blocks (e.g., individual transformer layers or weight projection blocks). For layer block $l \in \{0, 1, \dots, L-1\}$, the task vector is denoted as $\mathbf{\Delta}_{k, l} \in \mathbb{R}^{M_l}$, where $\sum_{l=0}^{L-1} M_l = M$.

In the standard Task Arithmetic baseline \cite{ilharco2022editing}, the task vectors are merged using a single uniform scalar coefficient $\lambda_k$ per task across all layers:
\begin{equation}
\Theta_{\text{merged}, l} = \Theta_{\text{base}, l} + \sum_{k=1}^K \lambda_k \mathbf{\Delta}_{k, l}
\end{equation}
This approach completely ignores the layer-wise functional differences across deep networks.

\subsection{Unconstrained Adaptive Merging and the Overfitting-Optimizer Paradox}
To overcome the limitations of uniform scaling, adaptive model merging methods like AdaMerging \cite{yang2024adamerging} optimize independent merging coefficients $\lambda_{k, l} \in \mathbb{R}$ for each task $k$ and layer block $l$. At test-time, the merged model weights are defined as:
\begin{equation}
\Theta_{\text{merged}, l}(\boldsymbol{\lambda}) = \Theta_{\text{base}, l} + \sum_{k=1}^K \lambda_{k, l} \mathbf{\Delta}_{k, l}
\end{equation}
where $\boldsymbol{\lambda} = \{ \lambda_{k, l} \} \in \mathbb{R}^{K \times L}$ represents the full set of learnable coefficients.

Because test-time adaptation is entirely unsupervised, $\boldsymbol{\lambda}$ is optimized on-the-fly using a surrogate objective over unlabeled test-time adaptation streams. Specifically, for each task $k$, we retrieve a batch of unlabeled inputs $x \sim \mathcal{D}_k^{\text{unlabeled}}$ and minimize the Shannon entropy of the model's predicted class distributions:
\begin{equation}
\label{eq:entropy_loss}
\mathcal{L}_{\text{TTA}}(\boldsymbol{\lambda}) = \sum_{k=1}^K \mathbb{E}_{x \sim \mathcal{D}_k^{\text{unlabeled}}} \left[ H\left(f_{\Theta_{\text{merged}}(\boldsymbol{\lambda})}(x)\right) \right]
\end{equation}
where $H(\mathbf{p}) = -\sum_{c=1}^C p_c \log (p_c + \epsilon)$ is the entropy of the probability vector $\mathbf{p} = f_{\Theta}(x)$ over $C$ target classes, and $\epsilon > 0$ is a microscopic numerical stability parameter (typically set to $10^{-8}$) to prevent taking the logarithm of zero when prediction probabilities approach extreme values during backpropagation.

\paragraph{The Overfitting-Optimizer Paradox.} Under unconstrained AdaMerging, the learnable parameter space has dimensionality $K \times L$ (or up to $K \times 52$ when optimizing projection-wise weights). Since the test-time adaptation streams are typically small and unlabeled, first-order optimizers like Adam easily minimize the unsupervised surrogate objective \eqref{eq:entropy_loss} by exploiting the high-frequency degrees of freedom in $\boldsymbol{\lambda}$. This results in a highly jagged, oscillating coefficient profile across layers (high roughness, or high total variation):
\begin{equation}
\mathcal{R}_{\text{TV}}(\boldsymbol{\lambda}) = \sum_{k=1}^K \sum_{l=1}^{L-1} (\lambda_{k, l} - \lambda_{k, l-1})^2 \gg 0
\end{equation}
This jaggedness represents the optimizer fitting transductive noise in the local adaptation batch, resulting in a fragile weight configuration that collapses on out-of-distribution or held-out test data. This is the core of the \textbf{Overfitting-Optimizer Paradox}: minimizing the surrogate loss on the adaptation stream damages true generalization performance.

\paragraph{The Degenerate Entropy Minimization Trap.}
Beyond simple generalization collapse, a major and notorious failure mode in unsupervised test-time adaptation is \emph{degenerate entropy minimization}. Because the Shannon entropy loss \eqref{eq:entropy_loss} is minimized to zero when the model outputs a delta-distribution (absolute confidence) for any class, a highly overparameterized unconstrained optimizer can easily find a degenerate, trivial solution. For instance, the optimizer can adjust the independent layer coefficients $\boldsymbol{\lambda}_k$ to break the functional representation of intermediate layers, turning the network into a constant predictor that outputs a single class with absolute certainty (confidence 100\%) for every input in the adaptation stream. This yields a surrogate loss of exactly zero, but results in a complete accuracy collapse (0\% accuracy) on actual test distributions.

Formally, let $C$ be the number of classes. A degenerate model $f_{\Theta_{\text{merged}}(\boldsymbol{\lambda})}$ outputs:
\begin{equation}
f_{\Theta_{\text{merged}}(\boldsymbol{\lambda})}(x)_c = \delta_{c, c^*}, \quad \forall x \in \mathcal{D}_k^{\text{unlabeled}}
\end{equation}
which achieves $\mathcal{L}_{\text{TTA}} = 0$. In physical neural networks, such degenerate state-fusing requires highly discontinuous, disjointed coefficient profiles across adjacent layers (e.g., highly suppressing the task vector in a specific early attention layer to nullify activations, while boosting it in a later feed-forward layer to saturate a single logit). These profiles exhibit extremely high total variation ($\mathcal{R}_{\text{TV}}(\boldsymbol{\lambda}) \gg 0$).

Crucially, our proposed \textbf{PolyMerge} framework mathematically pre-empts and eliminates these degenerate collapse states. By parameterizing coefficients as a smooth, low-degree polynomial of normalized depth, we restrict the learnable space to a low-frequency subspace. Since a low-degree polynomial (e.g., $d=1$ or $d=2$) cannot represent highly disjointed, localized layer-wise suppression or amplification, the optimization trajectory is physically unable to access the sharp, overfit regions of the parameter space where degenerate collapse occurs. Thus, the smooth continuous parameterization acts as a structural defense mechanism, forcing the model to minimize entropy only through genuine multi-task alignment. For extremely deep architectures (such as 32- or 80-layer networks) where high-degree global polynomials risk numerical instability and Runge's phenomenon, this continuous parameterization scales gracefully by employing B-splines and piecewise continuous local trajectories (see Appendix~\ref{sec:appendix_deeper} for a formal scaling analysis).

\subsection{PolyMerge: Polynomial Subspace Parameterization}
To resolve this paradox, we propose \textbf{PolyMerge}. Instead of allowing each layer block's coefficient to be optimized independently, we project the coefficient space onto a continuous, low-dimensional polynomial subspace of the layer depth. Specifically, we define the merging coefficient $\lambda_{k, l}$ for task $k$ and layer block $l$ as a polynomial of degree $d$:
\begin{equation}
\label{eq:poly_formulation}
\lambda_{k, l}(\boldsymbol{\alpha}) = \sum_{j=0}^d \alpha_{k, j} \cdot \left( \frac{l}{L-1} \right)^j
\end{equation}
where:
\begin{itemize}
    \item $\boldsymbol{\alpha} = \{ \alpha_{k, j} \} \in \mathbb{R}^{K \times (d+1)}$ represents the small set of learnable polynomial parameters.
    \item $\bar{l} = \frac{l}{L-1} \in [0, 1]$ is the normalized layer depth.
\end{itemize}

\paragraph{Physical Regularization via Normalized Depth.} Normalizing the layer index to the interval $[0, 1]$ is mathematically crucial:
1. \textbf{Numerical Stability:} It bounds the polynomial basis functions $l^j$ to prevent numerical overflow, which would otherwise occur for higher-order terms in deep models (e.g., $l=11$ or $l=51$).
2. \textbf{Scale Invariance:} It makes the polynomial parameterization invariant to the total number of layers $L$, allowing the same polynomial coefficients to represent identical continuous trajectories across architectures of varying depths.
3. \textbf{Well-Conditioned Optimization:} It restricts the domain of the Vandermonde matrix of the polynomial system to a compact interval, improving the conditioning of gradient updates.

\paragraph{Dimensionality Reduction.} By parameterizing coefficients via \eqref{eq:poly_formulation}, we reduce the number of learnable parameters per task from $L$ to $d+1$. For a typical Vision Transformer with $L=12$ layers (or $L=52$ projections) and a degree $d=2$ quadratic polynomial, PolyMerge reduces the learnable parameter space by \textbf{75.0\%} and \textbf{94.2\%}, respectively. This massive parameter reduction hard-constrains the search space, pruning high-frequency degrees of freedom and making it mathematically impossible for the optimizer to fit high-frequency transductive noise.

\begin{proposition}[Analytical Noise Filtering]
\label{prop:noise_filter}
Let $\mathbf{V} \in \mathbb{R}^{L \times (d+1)}$ represent the Vandermonde matrix mapping the polynomial parameters $\boldsymbol{\alpha}_k$ to the layer coefficients $\boldsymbol{\lambda}_k$: $\boldsymbol{\lambda}_k = \mathbf{V}\boldsymbol{\alpha}_k$. Let $\mathbf{P} = \mathbf{V}(\mathbf{V}^T \mathbf{V})^{-1} \mathbf{V}^T \in \mathbb{R}^{L \times L}$ be the orthogonal projection matrix onto the column space of $\mathbf{V}$. 

For any unconstrained coefficient vector corrupted by zero-mean white noise $\boldsymbol{\eta} \sim \mathcal{N}(\mathbf{0}, \sigma^2 \mathbf{I})$, the expected squared norm of the projected noise is reduced by a factor of $\frac{d+1}{L}$:
\begin{equation}
\mathbb{E}\left[ \|\mathbf{P} \boldsymbol{\eta}\|_2^2 \right] = \frac{d+1}{L} \cdot \mathbb{E}\left[ \|\boldsymbol{\eta}\|_2^2 \right]
\end{equation}
Furthermore, if the noise vector is a high-frequency alternating sequence $\eta_l = z \cdot (-1)^l$, then for $d \ll L$, the projection onto the low-frequency polynomial basis yields $\mathbf{P}\boldsymbol{\eta} \approx \mathbf{0}$, meaning high-frequency transductive noise is mathematically filtered out of the learnable subspace.
\end{proposition}

This proposition, proven in Appendix~\ref{sec:appendix_robustness}, demonstrates that the superiority of PolyMerge in a noisy test-time environment is a mathematically guaranteed property of low-frequency subspace projections. Rather than an accidental empirical finding, the "overfitting-prevention" is an analytical design feature of our parameterization.

\paragraph{Mitigating Boundary Oscillations (Runge's Phenomenon).} Although normalizing the layer depth to the compact interval $[0, 1]$ substantially improves Vandermonde matrix conditioning, executing global higher-degree polynomial interpolation ($d \ge 3$) on deep networks (e.g., $L \ge 32$) remains mathematically susceptible to Runge's phenomenon, where wild oscillations develop near the interval boundaries. To systematically bypass these boundary instabilities in deep architectures, we recommend transitioning from a single global polynomial to localized continuous piecewise splines or B-splines, which restrict representation adjustments to local compact domains. We provide a rigorous, detailed formulation and numerical analysis of this scaling strategy in Appendix~\ref{sec:appendix_deeper}.

\subsection{SplineMerge: Piecewise Continuous Splines}
\label{sec:splinemerge_method}
To reconcile the low-pass filtering benefits of PolyMerge with the realistic layer heterogeneity of deep networks (where adjacent structural blocks like attention and MLP layers exhibit sudden step-wise sensitivity transitions), we propose a piecewise-continuous extension termed \textbf{SplineMerge}. Instead of a single global polynomial trajectory over all $L$ layers, we partition the layer index range into $B$ disjoint block groups: $\mathcal{B}_1, \dots, \mathcal{B}_B$. Within each partition $\mathcal{B}_b$, we parameterize the coefficients local to that block as a low-degree polynomial:
\begin{equation}
\lambda_{k, l}(\boldsymbol{\alpha}) = \sum_{j=0}^d \alpha_{k, b, j} \cdot \left( \frac{l - l_{\text{start}}^{(b)}}{l_{\text{end}}^{(b)} - l_{\text{start}}^{(b)}} \right)^j, \quad \forall l \in \mathcal{B}_b
\end{equation}
where $l_{\text{start}}^{(b)}$ and $l_{\text{end}}^{(b)}$ denote the layer boundaries of block $b$. This piecewise formulation provides high localized flexibility to capture sharp transitions at block boundaries (e.g., between early, middle, and late layer groups), while still strictly limiting the parameter dimensions to $B(d+1) \ll L$. This maintains a robust low-pass filtering effect within each block group to reject transductive noise.

For scaling to much deeper models (such as LLaMA architectures with $32$, $40$, or $80$ layers), simple high-degree global polynomials are vulnerable to numerical instability and Runge's phenomenon. Under these deep regimes, SplineMerge with adaptive B-splines provides a structurally superior parameterization, enabling stable continuous local interpolation across dozens of layers as discussed in Appendix~\ref{sec:appendix_deeper}.

\subsection{Controlled Emulation and Calibration Framework}
To study the Overfitting-Optimizer Paradox under precise, reproducible, and highly controlled conditions, we design a continuous weight-merging simulation and optimization landscape emulator ($L=12$ layers). Rather than executing physical forward passes and backpropagation on real model checkpoints---which are highly sensitive to confounding stochastic variables, device-specific configurations, and private datasets---our emulator mathematically models the functional behavior of actual deep networks using hand-crafted analytical distance and loss functions. The optimal target profiles, baseline accuracies, and transductive noise are calibrated directly from empirical findings and classification statistics reported in prior Vision Transformer (ViT-B/32) model-merging publications, ensuring a high-fidelity optimization testbed. To ensure absolute scientific rigor and address simplifying assumptions, we formulate two distinct emulation regimes: the \emph{Stylized Convex Sandbox} (Model I) and the \emph{Physically Grounded Coupled Non-Convex Stress-Test} (Model II).

\subsubsection{Model I: Stylized Convex Sandbox}
Model I serves as our analytical baseline sandbox, employing a decoupled, convex quadratic loss landscape and simple alternating sign transductive noise. This setting is designed to allow clear physical intuition and tractability for our mathematical proofs.

\paragraph{Optimal Layer Importance Profiles.} For each task $k$, we define a continuous optimal layer-importance profile $\lambda^*_{k, l}$ representing the physical layer-wise trends of expert models fine-tuned on the respective datasets:
\begin{itemize}
    \item \textbf{MNIST:} $\lambda^*_{\text{MNIST}, l} = 0.5 - 0.25 \bar{l}$ (representing early-layer focus).
    \item \textbf{FashionMNIST:} $\lambda^*_{\text{F-MNIST}, l} = 0.2 + 0.35 \sin(\pi \bar{l})$ (representing a balanced mid-layer peak).
    \item \textbf{CIFAR-10:} $\lambda^*_{\text{CIFAR}, l} = 0.1 + 0.45 \bar{l}^2$ (representing late-layer specialized focus).
    \item \textbf{SVHN:} $\lambda^*_{\text{SVHN}, l} = 0.4 - 0.35 (\bar{l} - 0.5)^2$ (representing mid-layer strength).
\end{itemize}
where $\bar{l} = \frac{l}{L-1} \in [0, 1]$ is the normalized layer depth.

\paragraph{Transductive Overfitting Noise.} We model this transductive noise $\eta_{k, l}$ using high-frequency alternating oscillations across adjacent layers:
\begin{equation}
\eta_{k, l} = z_k \cdot (-1)^l, \quad z_k \sim \mathcal{N}(0, 0.12^2)
\end{equation}
This noise represents the local statistical perturbations that unconstrained first-order optimizers fit during test-time adaptation. The simulated training entropy loss is formulated as a decoupled quadratic distance to this perturbed target:
\begin{equation}
\label{eq:model1_loss}
\begin{aligned}
\mathcal{L}_{\text{TTA}}^{\text{I}}(\boldsymbol{\lambda}) = \sum_{k=1}^K \bigg[ &0.5 + \frac{5.0}{L} \sum_{l=0}^{L-1} \Big( \lambda_{k, l} \\
&- (\lambda^*_{k, l} + \eta_{k, l}) \Big)^2 \bigg]
\end{aligned}
\end{equation}

\paragraph{Generalization Accuracy.} Generalization performance depends on the alignment of the learned coefficients $\boldsymbol{\lambda}_k$ with the optimal profile $\lambda^*_{k, l}$ on the physical target distribution:
\begin{equation}
\text{Acc}_k^{\text{I}}(\boldsymbol{\lambda}_k) = \text{Base}_k + \delta_k \left( 1.0 - \frac{\mathcal{D}(\boldsymbol{\lambda}_k, \boldsymbol{\lambda}^*_k)}{\mathcal{D}(\mathbf{0.3}, \boldsymbol{\lambda}^*_k)} \right)
\end{equation}
where $\mathcal{D}(\mathbf{u}, \mathbf{v}) = \frac{1}{L} \sum_{l=0}^{L-1} (u_l - v_l)^2$, representing the mean squared distance. The baseline values $\text{Base}_k$ and the scaling sensitivity parameters $\delta_k$ are calibrated directly on empirical literature and summarized in Table~\ref{tab:calibration}.

\subsubsection{Model II: Physically Grounded Coupled Non-Convex Stress-Test}
To model real-world neural network landscapes more faithfully, Model II introduces highly complex layer functional interactions, a highly non-convex loss function with sub-optimal local minima, and multi-scale transductive noise.

\paragraph{Layer-wise Non-linear Sensitivity, Bottlenecks, and Coupling.}
In physical neural networks, layers do not operate in isolation; instead, weight merging and activation routing are highly non-linear and coupled processes. A minor deviation in a critical "bottleneck layer" (such as the self-attention query/key projection layers in intermediate transformer blocks) can disrupt the representation geometry and cause immediate catastrophic activation collapse. Conversely, late feed-forward projections or output projection layers are often highly redundant, possessing high resilience to coefficient fluctuations. Furthermore, adjacent layers exhibit strong functional coupling; a weight fluctuation in layer $l$ can often be compensated for by a coordinated adjustment in layer $l+1$, whereas uncoordinated, high-frequency oscillations across adjacent layers destroy functional compatibility.

To model these non-linear, coupled properties within our continuous emulator, we discard the simple decoupled Euclidean distance and define a layer sensitivity covariance matrix $\boldsymbol{\Sigma} \in \mathbb{R}^{L \times L}$. We set early layers to have low sensitivity ($s_l = 0.6$), middle layers moderate sensitivity ($s_l = 1.0$), and late layers high sensitivity ($s_l = 1.6$), modeling the empirical fact that deep layers in Vision Transformers are vastly more sensitive to weight merging. Adjacent layers are coupled with a positive correlation of $0.5$, yielding:
\begin{equation}
\Sigma_{i, j} = \sqrt{s_i s_j} \cdot 0.5^{|i-j|}
\end{equation}
The off-diagonal elements in $\boldsymbol{\Sigma}$ mathematically capture inter-layer coupling, representing the fact that coordinated smooth deviations are penalized far less than uncoordinated, high-frequency "jagged" oscillations. Generalization accuracy under Model II is evaluated using the Mahalanobis distance under $\boldsymbol{\Sigma}^{-1}$:
\begin{equation}
\label{eq:model2_accuracy}
\begin{aligned}
\text{Acc}_k^{\text{II}}(\boldsymbol{\lambda}_k) = &\text{Base}_k \\
&+ \delta_k \left( 1.0 - \frac{\mathbf{d}_k^T \boldsymbol{\Sigma}^{-1} \mathbf{d}_k}{\mathbf{d}_{0, k}^T \boldsymbol{\Sigma}^{-1} \mathbf{d}_{0, k}} \right)
\end{aligned}
\end{equation}
where $\mathbf{d}_k = \boldsymbol{\lambda}_k - \boldsymbol{\lambda}^*_k$ and $\mathbf{d}_{0, k} = \mathbf{0.3} - \boldsymbol{\lambda}^*_k$.
The inverse covariance $\boldsymbol{\Sigma}^{-1}$ acts as a quadratic form that heavily penalizes uncoordinated local fluctuations in bottleneck layers, mimicking the physical collapse of activation manifolds in deep networks.

\paragraph{Non-Convex Rastrigin Loss Landscape.} Rather than a simple convex quadratic loss, the test-time adaptation loss contains severe non-convexities with numerous sharp, sub-optimal local minima:
\begin{equation}
\label{eq:model2_loss}
\begin{aligned}
\mathcal{L}_{\text{TTA}}^{\text{II}}(\boldsymbol{\lambda}) = \sum_{k=1}^K \bigg[ &0.5 + 1.5 \cdot \mathbf{e}_k^T \boldsymbol{\Sigma}^{-1} \mathbf{e}_k \\
&+ 0.03 \sum_{l=0}^{L-1} \Big( 1 - \cos\big(10 \pi e_{k, l}\big) \Big) \bigg]
\end{aligned}
\end{equation}
where $\mathbf{e}_k = \boldsymbol{\lambda}_k - \mathbf{t}_k$, $e_{k, l}$ is the $l$-th element of $\mathbf{e}_k$, and $\mathbf{t}_k = \boldsymbol{\lambda}^*_k + \boldsymbol{\eta}_k$ represents the noisy target profile.

\paragraph{Multi-Scale Overfitting Noise.} We model transductive noise as a mixture of three noise components: alternating sign noise $\boldsymbol{\eta}_{k}^{\text{alt}}$, independent white Gaussian noise $\boldsymbol{\eta}_{k}^{\text{white}} \sim \mathcal{N}(\mathbf{0}, 0.08^2 \mathbf{I})$, and Brownian random-walk noise $\boldsymbol{\eta}_{k}^{\text{Brown}}$ (where $\eta_{k, l}^{\text{Brown}} = \eta_{k, l-1}^{\text{Brown}} + \epsilon_l$):
\begin{equation}
\boldsymbol{\eta}_k = 0.5 \boldsymbol{\eta}_{k}^{\text{alt}} + 0.3 \boldsymbol{\eta}_{k}^{\text{white}} + 0.2 \boldsymbol{\eta}_{k}^{\text{Brown}}
\end{equation}

\begin{table}[h]
\centering
\caption{Calibrated baseline and sensitivity parameters from CLIP literature.}
\label{tab:calibration}
\begin{small}
\begin{tabular}{lcc}
\toprule
Dataset & Baseline ($\text{Base}_k$) & Sensitivity ($\delta_k$) \\
\midrule
MNIST & $92.71\%$ & $1.5\%$ \\
FashionMNIST & $81.64\%$ & $4.0\%$ \\
CIFAR-10 & $90.17\%$ & $2.5\%$ \\
SVHN & $73.24\%$ & $5.5\%$ \\
\bottomrule
\end{tabular}
\end{small}
\end{table}

\subsection{Optimization and Backpropagation}
We initialize the polynomial coefficients $\boldsymbol{\alpha}$ such that the initial merging coefficients default to a standard uniform Task Arithmetic baseline (e.g., $\alpha_{k, 0} = 0.3$ and $\alpha_{k, j > 0} = 0.0$, making $\lambda_{k, l} = 0.3$ for all layers). This ensures that the optimization begins from a stable, well-performing weight configuration.

During test-time adaptation, we optimize $\boldsymbol{\alpha}$ to minimize the unsupervised simulated entropy loss. The gradients of the loss with respect to the learnable polynomial parameters $\alpha_{k, j}$ are propagated through the polynomial synthesis using the chain rule:
\begin{equation}
\frac{\partial \mathcal{L}_{\text{TTA}}}{\partial \alpha_{k, j}} = \sum_{l=0}^{L-1} \frac{\partial \mathcal{L}_{\text{TTA}}}{\partial \lambda_{k, l}} \cdot \frac{\partial \lambda_{k, l}}{\partial \alpha_{k, j}} = \sum_{l=0}^{L-1} \frac{\partial \mathcal{L}_{\text{TTA}}}{\partial \lambda_{k, l}} \cdot \left( \frac{l}{L-1} \right)^j
\end{equation}
This formulation allows us to easily compute gradients for any polynomial degree $d$ using standard auto-differentiation in PyTorch, maintaining extreme computational efficiency while enforcing global smoothness.
